% ============================================================
%  Anexo C: desarrollo estadístico e incertidumbre
%  Datos: evidencia_banco/datos/banco_pruebas_2026-11-17.db
% ============================================================
\section{Desarrollo estadístico y análisis de incertidumbre}
\label{anx:estadistica}

Este anexo desarrolla los cálculos cuyos resultados se presentan en la
Sección~\ref{subsec:resultados_consumo}. Se analizan dos series del soldador
registradas el 17/11/2026 en la base de datos del banco:
\begin{itemize}
    \item \textbf{Serie~A:} los cuatro ensayos de final de línea posteriores a la
          corrección de la estabilización (cada valor es el promedio del
          intervalo de ensayo).
    \item \textbf{Serie~B:} las ocho muestras del perfil de consumo tomadas con el
          contactor ya estabilizado (se excluyen las muestras del transitorio de
          cierre).
\end{itemize}

\begin{table}[htbp]
\centering
\caption{Valores de potencia activa de las series analizadas}
\label{tab:series}
\renewcommand{\arraystretch}{1.25}
\begin{tabular}{@{} l l @{}}
\toprule
\textbf{Serie} & \textbf{Potencia activa $P_i$ [\si{\watt}]} \\
\midrule
A (final de línea, $V \approx \SI{228}{\volt}$) & 36,8; 36,7; 36,7; 36,6 \\
B (perfil, $V \approx \SI{227}{\volt}$)
  & 36,1; 36,2; 36,1; 36,2; 36,3; 36,4; 36,4; 36,2 \\
\bottomrule
\end{tabular}
\end{table}

\subsection{Estadísticos descriptivos}

Para cada serie de $N$ valores se calculan la media, la desviación estándar
experimental, la incertidumbre tipo~A de la media y el coeficiente de variación:
\begin{align}
    \bar{P} &= \frac{1}{N}\sum_{i=1}^{N} P_i
    \label{eq:media} \\
    s_P &= \sqrt{\frac{1}{N-1}\sum_{i=1}^{N}\left(P_i - \bar{P}\right)^2}
    \label{eq:desvio} \\
    u_A &= \frac{s_P}{\sqrt{N}}
    \label{eq:tipoA} \\
    \mathrm{CV} &= \frac{s_P}{\bar{P}} \times \SI{100}{\percent}
    \label{eq:cv}
\end{align}

\subsection{Incertidumbre de medición}

Siguiendo la Guía para la Expresión de la Incertidumbre de Medición
\cite{gum}, se consideran como fuentes independientes:
\begin{itemize}
    \item \textbf{Tipo~A}, la dispersión de las mediciones, según la
          ecuación~\eqref{eq:tipoA}.
    \item \textbf{Tipo~B}, la exactitud declarada del medidor
          (\SI{0,5}{\percent} de la lectura), con distribución rectangular:
          \begin{equation}
              u_{B} = \frac{0{,}005\,\bar{P}}{\sqrt{3}}
              \label{eq:tipoB}
          \end{equation}
    \item \textbf{Tipo~B}, la resolución del medidor en potencia
          (\SI{0,1}{\watt}), con distribución rectangular:
          \begin{equation}
              u_{\mathrm{res}} = \frac{\SI{0,1}{\watt}}{2\sqrt{3}}
                               \approx \SI{0,029}{\watt}
              \label{eq:resolucion}
          \end{equation}
\end{itemize}
La incertidumbre combinada y la expandida, con factor de cobertura $k = 2$
(nivel de confianza aproximado del \SI{95}{\percent}), son:
\begin{align}
    u_c &= \sqrt{u_A^2 + u_B^2 + u_{\mathrm{res}}^2}
    \label{eq:combinada} \\
    U &= k \, u_c
    \label{eq:expandida}
\end{align}

La exactitud del medidor se toma de su hoja de datos; al no haberse contrastado
con un patrón, no se incluyen el error sistemático del conjunto (transformador de
corriente, cableado y conexiones) ni las correcciones de calibración. La variación
de la tensión de red durante cada serie fue menor a \SI{1}{\volt} y su efecto
queda contenido en la dispersión tipo~A.

\subsection{Resultados}

\begin{table}[htbp]
\centering
\caption{Estadísticos e incertidumbre de la potencia medida}
\label{tab:estadisticos}
\renewcommand{\arraystretch}{1.25}
\begin{tabular}{@{} l c c @{}}
\toprule
\textbf{Magnitud} & \textbf{Serie A} & \textbf{Serie B} \\
\midrule
$N$                                            & 4       & 8 \\
Media $\bar{P}$ [\si{\watt}]                   & 36,700  & 36,238 \\
Mínimo y máximo [\si{\watt}]                   & 36,6 -- 36,8 & 36,1 -- 36,4 \\
Desviación estándar $s_P$ [\si{\watt}]         & 0,082   & 0,119 \\
Coeficiente de variación [\si{\percent}]       & 0,22    & 0,33 \\
$u_A$ [\si{\watt}]                             & 0,041   & 0,042 \\
$u_B$ (exactitud) [\si{\watt}]                 & 0,106   & 0,105 \\
$u_{\mathrm{res}}$ [\si{\watt}]                & 0,029   & 0,029 \\
$u_c$ [\si{\watt}]                             & 0,117   & 0,116 \\
$U$ ($k = 2$) [\si{\watt}]                     & 0,23    & 0,23 \\
Desviación respecto de \SI{40}{\watt} [\si{\percent}] & $-8{,}3$ & $-9{,}4$ \\
\bottomrule
\end{tabular}
\end{table}

El resultado de la serie~A se expresa como
\begin{equation}
    P = (\num{36,70} \pm \num{0,23})~\si{\watt} \qquad (k = 2)
    \label{eq:resultado_final}
\end{equation}
La incertidumbre expandida representa el \SI{5,8}{\percent} de la semibanda de
tolerancia (\SI{4}{\watt}), por debajo del cociente de un cuarto que se suele
exigir entre la incertidumbre del instrumento y la tolerancia a verificar. La
contribución dominante es la exactitud declarada del medidor; la repetibilidad
aporta menos de la mitad.

\subsection{Efecto de la tensión de ensayo}

La cláusula~10.1 establece que el consumo se mide a la tensión nominal. En una
carga resistiva la potencia es proporcional al cuadrado de la tensión, por lo que
la potencia a la tensión nominal $V_n$ puede estimarse a partir de la medida a la
tensión $V$ como
\begin{equation}
    P_n = P \left(\frac{V_n}{V}\right)^2
    \label{eq:correccion_tension}
\end{equation}
El soldador está marcado para $V_n = \SI{230}{\volt}$. Con
$P = \SI{36,7}{\watt}$ a $V = \SI{227,9}{\volt}$ resulta
$P_n \approx \SI{37,4}{\watt}$, es decir, una desviación de \SI{-6,6}{\percent}
respecto del nominal: el aparato aprueba también a su tensión nominal y con mayor
margen. El efecto, sin embargo, no es despreciable: si el mismo aparato estuviera
marcado para \SI{220}{\volt}, la potencia estimada sería de \SI{34,2}{\watt}
(\SI{-14,5}{\percent}) y el veredicto sería FALLA. Para aparatos cercanos al
límite, el sistema debería alimentar el aparato a su tensión nominal o, en cargas
resistivas, aplicar esta corrección. Las lecturas del 15/11 a \SI{230,1}{\volt}
(\SIrange{37,5}{37,7}{\watt}) son coherentes con la
ecuación~\eqref{eq:correccion_tension}.
